\scanpage{8}
des corps vers les proobjets idoines, va avoir toutes les propriétés de fidélité du foncteur $K\unsure{-|\longrightarrow}B_K$, ou $K\unsure{-|\longrightarrow}\Pi_K$, ou $(K,\bar K)\unsure{-|\longrightarrow}\Gamma_{\bar K/K}$, \uline{restreint} aux corps \inserted{absolument} de type fini, auxquels nous allons par la suite nous borner, la plupart du temps. Mais il sera nécessaire, au cours du travail, de donner une description \inserted{purement} algébrique, p.ex., du pro-pro-groupe \unsure{fini} associé p.ex. à $\mathbb C$ (plus précisément, à $(\mathbb C,\mathbb C)$ !).

Le rôle dominant sera joué par le corps premier de car. $0$, $\mathbb Q$, donc par $B_{\mathbb Q}$ et $\Pi_{\mathbb Q}$, qui a un objet canonique, noté $\overline{\mathbb Q}_0$ -- la clôture alg. de $\mathbb Q$ dans $\mathbb C$. On pose
\[
(5)\qquad \Gamma_{\mathbb Q}=E_{\overline{\mathbb Q}_0/\mathbb Q}
\]
\margin{et on écrira \inserted{souvent} $\Gamma_{\overline{\mathbb Q}/\mathbb Q}$ au lieu de $E_{\overline{\mathbb Q}/\mathbb Q}$, pour une clôture alg. quelq. $\overline{\mathbb Q}$ de $\mathbb Q$}
Pour tt corps $K$ de car. $0$ -- en particulier pour les corps $K$ de type fini sur $\mathbb Q$, \unsure{aux-|au-}
\margin{7}
